\chapter[Requerimientos de red - Hardware]{Requerimientos de red - Hardware}
\label{cp:requerimientos-hardware}

\section*{Resumen}
Muchos de los laboratorios, además de necesitar bocas de red para las computadoras, necesitan bocas de red para otros equipos como proyectores, impresoras, equipos de medición, etc.

En esta parte del artículo solo se hará referencia a los requerimientos de hardware. 

\section{Necesidades de red}
\subsection{Ethernet}
Actualmente la infraestructura física está funcionando bien, con la cantidad de bocas necesarias para garantizar la conectividad y un desarrollo normal de las actividades.

\subsection{WiFi}
Actualmente los laboratorios tienen un access point cada uno marca TP-Link (\todo{buscar el modelo que ahora no lo recuerdo} ) instalados y al menos 2 de los 4 funcionando, el de Laboratorio 3 siempre tuvo problemas para conectar múltiples dispositivos y el de Laboratorio 1 también pero no tan frecuentemente.

Para el laboratorio 1, hace ya algunos años, se compraron access point nuevos marca Ubiquiti pero por desconocimiento del personal técnico del laboratorio, nunca se pudieron instalar. Por eso necesitamos que los instalen o que nos asesoren sobre cómo y donde instalarlos para garantizar una buena cobertura. También necesitamos que los configuren y mantengan.

\section{Impresoras}
Actualmente los laboratorios cuentan con:
\begin{itemize}
    \item 2 impresoras HP Laserjet HP4350
    \item 1 impresora HP Laserjet HP4350
    \item 1 impresora HP LaserJet Pro 400
    \item 1 impresora HP (\todo{Buscar modelo de la impresora A3 de tinta})
\end{itemize}

En un momento estuvieron conectadas a la red funcionando pero, actualmente solo 2 están conectadas. Si no recuerdo mal una de las impresoras funcionales, está conectada a la computadora del pañol por cable USB. 
Necesitamos que cada laboratorio tenga su impresora conectada a la red para que sea accesible por tecnicos y docentes.

Luego creo que con un mínimo asesoramiento el personal de los laboratorios será capaz de darles mantenimiento (reponer toner, hojas, destrabarlas en caso de atascos, etc.)

\section{Computadoras}
Actualmente y en lo que respecta al hardware el mantenimiento de las computadoras lo realiza el personal de los laboratorios sin mayores dificultades.
